% resume-default-medium.tex — the tagged resume under the default medium (#327).
%
% The body of resume.tex, whose wrappers pin `medium=print'; this one passes no
% `medium' option, so the class default, `screen', applies. The document is
% still two pages long, but neither page carries a folio and page two carries
% no running header, so page two holds no text artifact at all. The body's
% \href is underlined under `screen' (#278), and its anchor text must stay in
% the structure tree.
%
% The MuPDF and PDFKit baselines are resume.tex's with the three furniture
% lines removed. They were derived that way rather than written by --update, so
% `diff' against resume's shows the furniture and nothing else. The structure
% tree never held the furniture, so resume-default-medium.structure.txt is
% resume.structure.txt unchanged.
%
% The Poppler baseline differs in order as well. With the folio gone, each
% page's last entry is the last block Poppler sorts on its page, and the last
% entry on a page always trails its own dates and location column
% (tests/extraction/resume-entry-dates-order.tex, #219). So it records each
% page's dates after its bullets: expected output under that rule, written by
% --update and inspected against it (#327), not a defect.
\DocumentMetadata{lang=en,tagging=on}
\pdfvariable compresslevel=0
\pdfvariable objcompresslevel=0
\documentclass{careerdossier-resume}
\CDossierSetup{
  name={Tagged Industry Resume},
  headline={Accessibility Engineer},
  email={resume@example.test},
  phone={+1 555 0100},
  website={example.test/resume}
}
\begin{document}
\MakeCDossierHeader
\CDossierSection{Experience}
\begin{CDossierEntry}[
  title={Engineer},
  organization={Example Labs},
  location={Toronto},
  dates={2024--2026}
]
\begin{CDossierItemize}
  \item First resume achievement in source order.
  \item Second resume achievement with a meaningful
        \href{https://example.test/work}{work link}.
\end{CDossierItemize}
\end{CDossierEntry}
\newpage
\CDossierSection{Additional Experience}
\begin{CDossierEntry}[
  title={Senior Engineer},
  organization={Example Services},
  location={Ottawa},
  dates={2022--2024}
]
\begin{CDossierItemize}
  \item Second-page resume content after the shared page furniture.
\end{CDossierItemize}
\end{CDossierEntry}
\end{document}

